\chapter{Foundational Theorem Dossier}
\label{app:proetale-foundations}

This appendix is a ledger of the large theorems that support the proof.  Its
purpose is to prevent two opposite errors: pretending that the pro-etale site
makes finiteness, duality, or the trace formula formal, and treating the
Deligne--Laumon weight theorem as an unexplained black box.  Every entry has a
status, an exact interface, and a declared use.

\section{Status notation and global hypotheses}

The four statuses are:
\[
\begin{array}{c|l}
\mathbf F&
 \text{foundational theorem imported with an exact primary source},\\
\mathbf C&
 \text{construction or comparison carried out in the pro-etale chapters},\\
\mathbf P&
 \text{theorem proved in the main mathematical proof spine},\\
\mathbf A&
 \text{formal assembly from results already proved}.
\end{array}
\tag{A.1}\label{eq:dossier-statuses}
\]

Unless a statement says otherwise, schemes in the proof spine are
quasi-excellent, noetherian, and of finite type over a finite field;
morphisms are separated and of finite type.  In particular they are qcqs,
their structural maps are finitely presented, and \(\ell\) is invertible.
These hypotheses are stronger than necessary but simultaneously satisfy the
finiteness assumptions in the torsion and pro-etale six-operation theorems.

For the integral theory, \(E/\Ql\) is finite,
\(\mathcal O=\mathcal O_E\), \(\varpi\) is a uniformizer, and
\[
 \widehat{\mathcal O}_X
   =\varprojlim_n(\mathcal O/\varpi^n)_X
\]
is the sheaf associated with the topological ring \(\mathcal O\) on
\(X_{\proet}\).  The rational proof is performed over \(E\), followed by
finite coefficient extension when needed.

\section{Pro-etale topology and complete coefficients}

\begin{longtable}{@{}
 >{\raggedright\arraybackslash}p{0.12\textwidth}
 >{\raggedright\arraybackslash}p{0.24\textwidth}
 >{\raggedright\arraybackslash}p{0.39\textwidth}
 >{\raggedright\arraybackslash}p{0.15\textwidth}@{}}
\toprule
Status & Input & Exact interface & Use\\
\midrule
\endhead
\(\mathbf F\) &
Weakly etale topology &
A map is weakly etale when it and its diagonal are flat.  Weakly etale
objects with fpqc covering families define \(X_{\proet}\); etale maps are
exactly finitely presented weakly etale maps. &
Definition of the site and functoriality.\\

\(\mathbf F\) &
Weakly contractible basis &
Every scheme has a pro-etale cover by affine \(w\)-contractible objects.
Evaluation on such an object is exact and commutes with limits. &
Derived completion and local perfectness.\\

\(\mathbf F\) &
Repleteness &
For a countable sequential tower
\(F_0\leftarrow F_1\leftarrow\cdots\) of sheaves on the pro-etale topos in
which every transition is an epimorphism,
\(\varprojlim F_n\simeq\Rlim F_n\).  General countable sequential inverse
limits have cohomological dimension at most one, and the derived category is
left complete. &
Convergence of coefficient towers.\\

\(\mathbf F\) &
Etale comparison &
For torsion coefficients, pullback along
\(\nu:X_{\proet}\to X_{\et}\) is fully faithful on the bounded-below derived
category and computes the same cohomology. &
Transport of finite-level theorems.\\

\(\mathbf C\) &
Complete coefficient sheaf &
The sheaf associated with the topological ring \(\mathcal O\) is
\(\widehat{\mathcal O}_X=\varprojlim_n\mathcal O/\varpi^n\), and it is
derived complete. &
An honest integral coefficient object.\\

\(\mathbf C\) &
Constructible category &
\(K\in\Dcons(X_{\proet},\widehat{\mathcal O})\) exactly when \(K\) is
derived complete and \(K\otimes^L\mathcal O/\varpi\) is constructible of
finite Tor amplitude.  On topologically noetherian \(X\), this is equivalent
to local perfectness along a finite stratification. &
Primary derived category of the notes.\\

\(\mathbf C\) &
Coefficient change &
\[
 \begin{gathered}
  \Dcons(X,\widehat{\mathcal O})[\varpi^{-1}]\\[-0.2em]
  {}\simeq \Dcons(X,E),\\
  K\simeq\Rlim_nK/\varpi^n.
 \end{gathered}
\]
Every \(\Qbarl\)-constructible object and every finite diagram involving it
descends to some finite \(E/\Ql\). &
Lattices and rational weights.\\
\bottomrule
\end{longtable}

The primary source for the first six rows is
\href{https://people.mpim-bonn.mpg.de/scholze/proetale.pdf}
{Bhatt--Scholze, \emph{The pro-etale topology for schemes}},
especially Definition 1.2, Theorem 1.5, Proposition 3.1.10,
Corollary 5.1.6, Definition 6.5.1, Proposition 6.6.11, and
Proposition 6.8.14.  The characterization
\emph{lisse means dualizable, constructible means lisse on a
stratification} is developed systematically in
\href{https://arxiv.org/abs/2305.18131}
{Hemo--Richarz--Scholbach, \emph{Constructible sheaves on schemes}}.

\begin{warning}
Geometric points are conservative for the constructible category by
finite-level comparison and stratification.  The notes do not require, and
do not assert, that the usual set of geometric points detects every object
of the entire unbounded pro-etale derived category.
\end{warning}

\section{The six operations}

There are two logically separate claims: the functors exist on an ambient or
constructible derived category, and they preserve bounded constructibility.
The second contains the deep geometric input.

\begin{foundation}[title={Active architecture: Route B}]
The ambient category in this edition is
\(\Dsolid(X,\Lambda)\), the reflective localization of solid pro-etale
complexes.  It has the six operations before constructibility is imposed,
by solid pro-etale motives and rigidity
\cite[Theorem~3.3.7]{RuimyTubachWolfSolid}.  Constructible complexes then
embed as a full subcategory.  The table below records the proof that the
ambient operations preserve that subcategory and their comparison with the
classical adic theory; it is not a levelwise definition of the ambient
functors.
\end{foundation}

\begin{longtable}{@{}
 >{\raggedright\arraybackslash}p{0.11\textwidth}
 >{\raggedright\arraybackslash}p{0.20\textwidth}
 >{\raggedright\arraybackslash}p{0.45\textwidth}
 >{\raggedright\arraybackslash}p{0.14\textwidth}@{}}
\toprule
Status & Theorem & Form used & Dependency\\
\midrule
\endhead
\(\mathbf F\) &
Torsion finiteness &
For \(\Lambda_n=\mathcal O/\varpi^n\), finite-type direct images of bounded
constructible complexes have bounded constructible cohomology; the
cohomological amplitude is uniformly bounded in \(n\). &
SGA 4 and SGA \(4\frac12\).\\

\(\mathbf C\) &
Completed pullback &
The left adjoint on derived-complete objects is
\[
 \begin{aligned}
  Lf_{\comp}^*K&=\bigl(Lf^*K\bigr)^\wedge,\\
  (Lf_{\comp}^*K)/\varpi^n
    &\simeq Lf^*(K/\varpi^n).
 \end{aligned}
\]
It preserves constructibility. &
Completion plus torsion comparison.\\

\(\mathbf C\) &
Direct image &
\(Rf_*\) preserves derived completeness.  Under the stated finiteness
hypotheses it preserves constructibility and reduces to finite-level
\(Rf_*\) modulo every \(\varpi^n\). &
Uniform torsion finiteness.\\

\(\mathbf C\) &
Exceptional direct image &
For \(f=\bar f j\), with \(j\) open and \(\bar f\) proper, set
\[
 Rf_!K=R\bar f_*j_!K.
\]
Common refinements give canonical comparison isomorphisms satisfying the
cocycle condition. &
Nagata compactification and proper base change.\\

\(\mathbf C\) &
Exceptional pullback &
\(Rf^!\) is the right adjoint of \(Rf_!\) on the constructible category and
\[
 (Rf^!K)/\varpi^n\simeq Rf_n^!(K/\varpi^n).
\]
For smooth \(f\) of pure relative dimension \(d\),
\(Rf^!K\simeq Lf^*K(d)[2d]\). &
Finite-level duality and purity.\\

\(\mathbf C\) &
Tensor and internal Hom &
Constructible objects are closed under derived tensor and internal Hom.
The integral tensor is already derived complete on constructible objects;
both operations commute with finite reduction in their finite-Tor range. &
Local perfectness.\\

\(\mathbf F/\mathbf C\) &
Base change and projection formula &
Proper base change for \(Rf_*\), arbitrary base change for \(Rf_!\), smooth
base change, projection formula, localization, and Kunneth hold at finite
level and lift conservatively through all reductions. &
SGA finite-level theorems, then completion.\\
\bottomrule
\end{longtable}

The active ambient construction is
\cite[Theorems~2.5.4 and~3.3.7]{RuimyTubachWolfSolid}.  Its proof of proper
base change uses torsion constructible generators, so the torsion row of the
table remains a genuine foundational dependency.  Constructible preservation
and comparison with classical adic complexes are in Bhatt--Scholze
Section~6.7 and Hemo--Richarz--Scholbach Theorem~7.7.  The particularly
important nonexample remains Bhatt--Scholze Remark~6.7.8: \(Rf_!\) is not, in
general, the derived functor of a degree-zero abelian \(f_!\).  Route B uses
the coherent ambient compactification construction and never that invalid
shortcut.

\section{Duality, purity, and trace}

\begin{longtable}{@{}
 >{\raggedright\arraybackslash}p{0.11\textwidth}
 >{\raggedright\arraybackslash}p{0.21\textwidth}
 >{\raggedright\arraybackslash}p{0.45\textwidth}
 >{\raggedright\arraybackslash}p{0.13\textwidth}@{}}
\toprule
Status & Theorem & Exact form consumed & Source\\
\midrule
\endhead
\(\mathbf F\) &
Dualizing complex &
For excellent \(\ell\)-coprime \(X\) with a dimension function, there is
\(\omega_X\in\Dcons(X,\widehat{\mathcal O})\) such that
\[
 \D_XK=R\mathcal Hom(K,\omega_X),\qquad
 K\xrightarrow{\sim}\D_X\D_XK
\]
on constructible objects. &
ILO; Bhatt--Scholze 6.7.20.\\

\(\mathbf F\) &
Absolute cohomological purity &
For a regular closed immersion \(i:Z\hookrightarrow X\) of codimension
\(c\), with regular source and target and \(\ell\) invertible,
\[
 Ri^!\widehat{\mathcal O}_X
   \simeq\widehat{\mathcal O}_Z(-c)[-2c].
\]
&
Gabber purity.\\

\(\mathbf F/\mathbf C\) &
Verdier duality &
\[
 \D_YRf_!\simeq Rf_*\D_X,\qquad Rf_!\dashv Rf^!.
\]
compatible with composition, localization, and coefficient reduction. &
Finite level, then completion.\\

\(\mathbf F/\mathbf C\) &
Perfectness &
For \(X\) separated of finite type over a separably closed or finite field
and \(K\) constructible of finite Tor amplitude,
\(R\Gamma_c(X,K)\) is a perfect coefficient complex. &
Finiteness and topological Nakayama.\\

\(\mathbf F\) &
Lefschetz--Verdier trace &
For a constructible Weil complex and a proper fixed-point correspondence,
the categorical trace equals the sum of local terms.  For geometric
Frobenius on a finite-type scheme,
\[
 \begin{aligned}
 &\sum_{x\in X_0(\mathbf F_{q^n})}
   \Tr(F_x^n\mid K_{\bar x})\\[-0.2em]
 &\qquad =\sum_i(-1)^i
   \Tr(F^n\mid H_c^i(X,K)).
 \end{aligned}
\]
&
SGA 5 and SGA \(4\frac12\).\\
\bottomrule
\end{longtable}

In the rational-point formula, \(F_x^n\) means the action at the chosen
\(\mathbf F_{q^n}\)-point.  If it is rewritten as a sum over closed points
of degree \(e\mid n\), the operator is \(F_x^{\,n/e}\).  This conversion is
recorded separately so that no exponent convention is hidden.

\section{Perverse foundations and the integral warning}

The middle perverse \(t\)-structure over \(E\) or \(\Qbarl\) consumes the
following non-arithmetic statements from BBD \cite{BBD}:

\begin{longtable}{@{}
 >{\raggedright\arraybackslash}p{0.12\textwidth}
 >{\raggedright\arraybackslash}p{0.31\textwidth}
 >{\raggedright\arraybackslash}p{0.31\textwidth}
 >{\raggedright\arraybackslash}p{0.16\textwidth}@{}}
\toprule
Status & Input & Form used & Primary location\\
\midrule
\endhead
\(\mathbf F/\mathbf P\) &
Gluing of the middle \(t\)-structure &
Support and cosupport conditions glue over a finite smooth
stratification, are independent of refinement, and give a bounded
\(t\)-structure. &
BBD 1.4, 2.1--2.2.\\

\(\mathbf F/\mathbf P\) &
Affine Artin vanishing &
For affine \(f\), \(Rf_*\) is right perverse \(t\)-exact and \(Rf_!\) is
left perverse \(t\)-exact; quasi-finite affine maps are \(t\)-exact. &
BBD 4.1.1--4.1.3.\\

\(\mathbf F/\mathbf P\) &
Functorial amplitudes &
Smooth pullback shifted by relative dimension, finite direct image, and
external product are perverse \(t\)-exact. &
BBD 4.2.4, 4.2.7--4.2.8.\\

\(\mathbf P\) &
Intermediate extension &
\(j_{!*}P=\im({}^pj_!P\to{}^pj_*P)\), characterized by the absence of
boundary subobjects and quotients. &
BBD 1.4.23 and 4.3.1.\\

\(\mathbf F/\mathbf P\) &
Finite length and simple objects &
The rational perverse heart has finite length; its simple objects are
intersection complexes of irreducible lisse sheaves on smooth strata. &
BBD 4.3.1.\\
\bottomrule
\end{longtable}

Over \(\mathcal O\), Verdier duality does not preserve the ordinary
\(p\)-perverse heart.  It exchanges \(p\) and \(p^+\), the tilt by the
torsion/torsion-free pair.  This is BBD Section 3.3.  The notes prove the
comparison with rationalization and use the self-dual middle perversity only
after passage to \(E\).  No integral decomposition theorem is assumed.

\section{Arithmetic inputs upstream of the main weight theorem}

The following inputs are permitted before the general direct-image theorem.
None is smooth-proper purity in arbitrary dimension.

\begin{longtable}{@{}
 >{\raggedright\arraybackslash}p{0.11\textwidth}
 >{\raggedright\arraybackslash}p{0.24\textwidth}
 >{\raggedright\arraybackslash}p{0.43\textwidth}
 >{\raggedright\arraybackslash}p{0.12\textwidth}@{}}
\toprule
Status & Input & Exact role & Source\\
\midrule
\endhead
\(\mathbf F\) &
Rank-one theory &
Class field theory and determinant arguments control weights for rank-one
lisse Weil sheaves. &
Weil II 1.2--1.4.\\

\(\mathbf F\) &
Local monodromy weights &
For a pure lisse sheaf on a punctured trait, the graded pieces of the
monodromy filtration have the prescribed shifted weights. &
Weil II 1.8.4.\\

\(\mathbf F/\mathbf P\) &
Real-sheaf criterion &
For an \(\iota\)-real lisse Weil sheaf, every irreducible constituent is
\(\iota\)-pure, with the compatible pairing of weights in Deligne's
criterion.  A finite-order determinant after a Tate or constant twist is
used to normalize the determinant and the weight; that condition alone does
not imply purity. &
Weil II 1.5.1.\\

\(\mathbf F\) &
Fourier support-forgetting &
For the Artin--Schreier Fourier kernel, the canonical map from
\(!\)-pushforward to \(*\)-pushforward is an isomorphism. &
Laumon 1.3.1.1.\\

\(\mathbf P\) &
Fourier formalism &
Inversion, duality with the inverse character, Plancherel, and perverse
\(t\)-exactness are derived from the six operations, character
orthogonality, support-forgetting, and affine exactness. &
Laumon 1.2--1.3.\\

\(\mathbf P\) &
Curve estimate &
For lisse input of weights at most \(w\), the relevant
\(H_c^1\) eigenvalues have weights at most \(w+1\), with local boundary
terms controlled separately. &
Laumon 4; Weil II 3.2.3.\\
\bottomrule
\end{longtable}

The Artin--Schreier kernel is defined over a finite \(E/\Ql\) containing the
values of the chosen additive character and has an
\(\mathcal O_E\)-lattice.  It is not silently assumed to be defined over
\(\Zl\).

\section{Outputs proved in the notes}

The following statements are not foundational imports:

\begin{enumerate}
\item the perverse Fourier mechanism and the curve weight estimate;
\item the perverse-first theorem
      \[
       \operatorname{WeightsLE}(\cP,w)
       \Longrightarrow
       \operatorname{WeightsLE}(\pH^aRf_!\cP,w+a);
      \]
\item the derived \(Rf_!\) upper-weight theorem obtained from the perverse
      truncation tower;
\item the dual \(Rf_*\) lower bound and proper direct-image purity;
\item purity of \(j_{!*}\), intersection complexes, and proper intersection
      cohomology;
\item geometric semisimplicity of pure perverse sheaves and the
      decomposition theorem after geometric base change;
\item purity of \(H^i\) for smooth proper varieties;
\item rationality, the functional equation, integrality of the individual
      factors, and independence of \(\ell\).
\end{enumerate}

\section{Terminal dependency audit}

\[
\begin{tikzcd}[column sep=small,row sep=small]
\substack{\mathbf C:\ \text{solidification and}\\
                 \text{pro-etale motives}}
 \arrow[r] &
\substack{\mathbf C:\ \text{solid rigidity}}
 \arrow[r] &
\substack{\mathbf C:\ \text{ambient coherent}\\
                 \text{six operations}} \\
\substack{\mathbf F:\ \text{torsion base change}\\
                 \text{and rigidity}}
 \arrow[ur] \arrow[r] &
\substack{\mathbf F:\ \text{finiteness, duality,}\\
                 \text{purity, and trace}}
 \arrow[r] &
\substack{\mathbf C:\ \text{restriction to}\\
                 \Dcons}
 \arrow[u] \arrow[d] \\
\substack{\mathbf A:\ \text{classical}\\
                 \text{Weil conjectures}}
 &
\substack{\mathbf P:\ \text{proper purity,}\\
                 \text{IC and decomposition}}
 \arrow[l] &
\substack{\mathbf P:\ \text{perverse Fourier}\\
                 \text{weight engine}}
 \arrow[l]
\end{tikzcd}
\tag{A.2}\label{eq:dossier-terminal-dependencies}
\]

There is no reverse arrow from smooth-proper purity, decomposition, hard
Lefschetz, integral Frobenius factors, or \(\ell\)-independence to the
Fourier or direct-image nodes.
